% TOP_PROBLEM: {"problemId": "problem.voronoi-conjecture-on-parallelohedra", "canonicalTitle": "Voronoi Conjecture on Parallelohedra", "releaseRank": 477, "primaryDomain": "combinatorics_discrete_geometry", "primaryDomainLabel": "Combinatorics & discrete geometry", "displayStatus": "open_with_solved_subcases", "releaseStatus": "open_with_solved_subcases"}
% !TeX program = lualatex
% Definition and sources only; this file is not an LLM solution attempt.
\documentclass[11pt]{article}
\usepackage[margin=25mm]{geometry}
\usepackage{fontspec}
\setmainfont{Latin Modern Roman}
\usepackage{amsmath,amssymb,mathtools}
\usepackage{unicode-math}
\setmathfont{Latin Modern Math}
\usepackage{xurl,hyperref}
\hypersetup{colorlinks=true,urlcolor=blue}
\setcounter{secnumdepth}{0}
\setlength{\parindent}{0pt}
\setlength{\parskip}{5pt}
\setlength{\emergencystretch}{3em}
\begin{document}
\section*{\texorpdfstring{Voronoi Conjecture on Parallelohedra}{Voronoi Conjecture on Parallelohedra}}\label{477}
\subsection{Definitions and mathematical statement}
A parallelohedron is a full-dimensional convex polytope $P\subset{\mathbb{R}}^d$ that tiles ${\mathbb{R}}^d$ by translations, with disjoint interiors of distinct tiles. For a full-rank lattice $\Lambda$, its Euclidean Dirichlet--Voronoi cell at zero is
\[
 V(\Lambda)=\{x\in{\mathbb{R}}^d:\|x\|\le\|x-v\|
                 \text{ for every }v\in\Lambda\}.
\]
The conjecture asks whether every parallelohedron is of the form $A V(\Lambda)+b$ for some invertible linear map $A$, vector $b$, and lattice $\Lambda$. An equivalent viewpoint allows a positive-definite quadratic form in place of the Euclidean norm. The requested equivalence is affine, not necessarily an isometry or a similarity. Finding a translational tiling alone does not identify it as a Voronoi cell after a linear change of coordinates.
\par\smallskip\textit{Formulation and concept references:} \href{https://arxiv.org/abs/1906.05193}{[S1]}.

\subsection{Short English statement}
Can every convex polytope that tiles space by translation be transformed affinely into a lattice's nearest-point region?
\subsection{Sources}
\begin{itemize}
\item[S1]A. Garber, Voronoi conjecture for five-dimensional parallelohedra, Inventiones Mathematicae (2025); arXiv:1906.05193.
\url{https://arxiv.org/abs/1906.05193}.
\textit{Formulation source.}
\end{itemize}
\end{document}
